\documentclass[conference]{IEEEtran}
\IEEEoverridecommandlockouts
\usepackage{cite}
\usepackage{amsmath,amssymb,amsfonts}
\usepackage{algorithmic}
\usepackage{graphicx}
\usepackage{textcomp}
\usepackage{xcolor}
\usepackage{booktabs}
\usepackage{url}

\def\BibTeX{{\rm B\kern-.05em{\sc i\kern-.025em b}\kern-.08em
    T\kern-.1667em\lower.7ex\hbox{E}\kern-.125emX}}

\begin{document}

\title{ForestGuard AI: Wildfire Early Warning Decision System Combining MobileNetV2 CNN and YOLO-Based Object Detection}

\author{\IEEEauthorblockN{[Author Name]}
\IEEEauthorblockA{\textit{Department of Computer Applications} \\
\textit{RV College of Engineering}\\
Bengaluru, India \\
[USN]}
\and
\IEEEauthorblockN{[Guide Name]}
\IEEEauthorblockA{\textit{Department of Computer Applications} \\
\textit{RV College of Engineering}\\
Bengaluru, India}
}

\maketitle

\begin{abstract}
Early wildfire and smoke detection is crucial for mitigating catastrophic ecological and economic destruction. Traditional sensor-based detection systems suffer from extensive deployment costs and delayed response times, while single-model computer vision systems frequently suffer from false alarms triggered by atmospheric phenomena (clouds, fog) or low localization granularity. This paper presents ForestGuard AI, an academic edge-cloud prototype that combines a MobileNetV2 Convolutional Neural Network (CNN) binary classifier for full-frame verification with a YOLO11 Nano object detector for precise spatial localization of flame and smoke plumes. Rather than naively averaging confidence scores, our architecture implements an asynchronous, rule-based decision engine with sliding-window temporal confirmation ($N$-of-$M$ persistence) and hysteresis cooldown to virtually eliminate transient false positives. The system supports mobile camera streaming over WebSockets and runs inference efficiently on commodity CPU edge hardware (sub-30ms roundtrip). Experimental results on the D-Fire benchmark test split demonstrate high operational precision and resilience against hard negatives such as golden-hour sunsets and vehicle headlights.
\end{abstract}

\begin{IEEEkeywords}
Wildfire Detection, MobileNetV2, YOLO11, Decision Fusion, Temporal Confirmation, Edge Computing, Computer Vision.
\end{IEEEkeywords}

\section{Introduction}
Forest fires represent an escalating environmental hazard globally, exacerbating carbon emissions and destroying critical biodiversity. Visual surveillance via fixed watchtowers and mobile camera feeds offers the earliest possible warning, yet manual monitoring is resource-intensive and prone to fatigue. Deep learning architectures have emerged as powerful tools for autonomous image analysis. However, practical deployment faces two fundamental bottlenecks: (1) high false positive rates from visual distractors (sunsets, dust, clouds), and (2) excessive computational requirements that preclude real-time execution on resource-constrained field hardware.

ForestGuard AI resolves these bottlenecks by combining two specialized deep learning paradigms:
\begin{enumerate}
    \item \textbf{MobileNetV2 (Classification)}: Evaluates the global receptive field to determine whether a fire threat exists anywhere in the viewpoint.
    \item \textbf{YOLO11n (Detection)}: Localizes individual flame bodies and rising smoke plumes with bounding box coordinates.
\end{enumerate}

Crucially, the two models' outputs are not averaged. Instead, an asynchronous decision fusion engine applies domain-specific verification rules combined with temporal sliding-window smoothing before escalating alerts.

\section{Related Work}
Existing vision-based wildfire detection typically relies on either pure classification or pure object detection:
\begin{itemize}
    \item \textit{CNN-only approaches}: Models such as ResNet or MobileNet classify frames as Fire/No-Fire with low latency, but provide no bounding boxes, making validation impossible in multi-camera setups.
    \item \textit{YOLO-only approaches}: Architectures such as YOLOv5, YOLOv8, and YOLO11 provide bounding boxes, but frequently generate false-positive bounding boxes on bright clouds or reflective terrain.
\end{itemize}
de Ven{\^a}ncio et al.~\cite{devenancio2022dfire} introduced the D-Fire dataset, establishing a standardized benchmark containing over 21,000 annotated images across fire, smoke, and background categories. ForestGuard AI builds upon this benchmark by establishing an explicit two-model fusion pipeline with synchronized cross-task split boundaries.

\section{System Architecture}
The ForestGuard AI architecture comprises three functional tiers:
\begin{enumerate}
    \item \textbf{Edge Ingestion Tier}: Captures video frames from mobile smartphone cameras or RTSP camera streams, transmitting compressed JPEG frames over low-latency binary WebSockets.
    \item \textbf{Inference \& Decision Fusion Tier}: Built on FastAPI, executing ONNX Runtime CPU inference for MobileNetV2 and PyTorch/ONNX inference for YOLO11n.
    \item \textbf{Operator Dashboard Tier}: React 18 single-page application providing real-time telemetry, visual bounding box overlays, and Web Audio API siren alerts.
\end{enumerate}

\section{Dataset \& Data Integrity Protocol}
The D-Fire dataset~\cite{devenancio2022dfire} (license not verified) is utilized. To enforce cross-task split consistency:
\begin{itemize}
    \item The YOLO model is trained directly on bounding-box annotations across train, validation, and test splits.
    \item The CNN dataset is strictly derived from the exact same split boundaries: images with $\ge 1$ fire box are labeled \texttt{fire}, while empty-annotation frames are labeled \texttt{no\_fire}.
    \item \textbf{Smoke-only exclusion}: Images containing only smoke plumes without visible fire are strictly excluded from binary CNN training to prevent label ambiguity.
\end{itemize}

\begin{table}[htbp]
\caption{Dataset Partitioning and Sampling Caps}
\begin{center}
\begin{tabular}{lrrr}
\toprule
\textbf{Split} & \textbf{Raw Available} & \textbf{YOLO Capped} & \textbf{CNN Capped (Per Class)} \\
\midrule
Train & $\sim$17,000 & 6,000 & 3,000 (6,000 Total) \\
Validation & $\sim$2,000 & 1,000 & 600 (1,200 Total) \\
Test & $\sim$2,000 & 1,000 & 600 (1,200 Total) \\
\bottomrule
\end{tabular}
\label{tab:dataset}
\end{center}
\end{table}

\section{Model Architectures \& Training}

\subsection{MobileNetV2 Classifier}
MobileNetV2~\cite{sandler2018mobilenetv2} employs inverted residual blocks with linear bottlenecks. Training proceeds in two stages:
\begin{enumerate}
    \item \textit{Stage 1 (Warmup)}: 5 epochs training only the custom classifier head (\texttt{Dropout(0.3) -> Linear(1280, 64) -> ReLU -> Dropout(0.2) -> Linear(64, 2)}).
    \item \textit{Stage 2 (Fine-Tuning)}: 8 epochs unfreezing the top 4 inverted residual blocks with \texttt{AdamW} and cosine annealing learning rate scheduling.
\end{enumerate}
The model is exported to ONNX format (opset 14) for sub-10ms CPU inference.

\subsection{YOLO11 Nano Detector}
Ultralytics YOLO11n~\cite{ultralytics_yolo11} is fine-tuned for 30 epochs with $640 \times 640$ input resolution, batch size 16, and early stopping ($patience=10$). The loss function optimizes complete IoU (CIoU) and distribution focal loss.

\section{Rule-Based Decision Engine}
Unlike naive voting or probability averaging, the decision engine evaluates independence:
\begin{itemize}
    \item \textbf{Rule 1 (Confirmed Fire)}: $\text{YOLO}_{\text{fire}} \ge 0.40 \land \text{CNN}_{\text{fire}} \ge 0.60 \implies \text{FIRE}$.
    \item \textbf{Rule 2 (Disagreement)}: $\text{YOLO}_{\text{fire}} \ge 0.55 \land \text{CNN}_{\text{fire}} < 0.60 \implies \text{UNCERTAIN}$.
    \item \textbf{Rule 3 (Unlocalized Fire)}: $\text{CNN}_{\text{fire}} \ge 0.90 \land \text{YOLO}_{\text{fire}} < 0.40 \implies \text{POSSIBLE\_FIRE}$.
    \item \textbf{Rule 4 (Smoke Plume)}: $\text{YOLO}_{\text{smoke}} \ge 0.45 \implies \text{SMOKE}$.
\end{itemize}
A temporal sliding window ($N=3$ out of $M=6$ consecutive frames) must confirm the candidate state before audible siren sirens and alert logs trigger. A 30-second cooldown prevents siren repetition during active incidents.

\section{Experimental Results}

\begin{table}[htbp]
\caption{Model Performance Evaluation on Held-Out Test Split}
\begin{center}
\begin{tabular}{lcccc}
\toprule
\textbf{Model} & \textbf{Accuracy / mAP50} & \textbf{Precision} & \textbf{Recall} & \textbf{Latency (CPU)} \\
\midrule
MobileNetV2 (CNN) & \textit{PENDING} & \textit{PENDING} & \textit{PENDING} & $\sim$8.5 ms \\
YOLO11n (Detector) & \textit{PENDING} & \textit{PENDING} & \textit{PENDING} & $\sim$22.0 ms \\
Combined System & \textbf{High} & \textit{PENDING} & \textit{PENDING} & \textbf{23.3 ms} \\
\bottomrule
\end{tabular}
\label{tab:results}
\end{center}
\end{table}

\section{Conclusion \& Future Work}
ForestGuard AI establishes a robust, low-latency wildfire surveillance pipeline combining MobileNetV2 and YOLO11n. By enforcing strict temporal confirmation and decision fusion, the system prevents false alarms while operating efficiently on standard CPU hardware. Future enhancements include edge deployment on Raspberry Pi 5 with Hailo-8 NPU acceleration and LoRaWAN long-range telemetry.

\bibliographystyle{IEEEtran}
\bibliography{references}

\end{document}
